\chapter{Communication Complexity — Hard Sample}
%%% generated by the blueprint pipeline from TCSlib.CommunicationComplexity.NewmanTheorem.FuncDisjointnessLowerBound.HardSample
\section{Overview}

This module builds the explicit finite sample space carrying Razborov's hard distribution
for set disjointness, used in the information-complexity proof of the linear randomized
lower bound: a uniformly random special coordinate $T$, two independent uniform bits at
$T$, and at every other coordinate a uniformly random pair from $\{00, 01, 10\}$.  It also
provides the uniform laws on the factors, the accessors reading off Alice's and Bob's
inputs, and the basic events and conditioning variables used downstream.

%%%%%%%%%%%%%%%%%%%%%%%%%%%%%%%%%%%%%%%%%%%%%%%%%%%%%%%%%%%%%%%%%%%%%%%%%%%%%%%
\section{Declarations}

\begin{definition}[Disjoint coordinate pair]
\lean{CommunicationComplexity.Functions.Disjointness.RandomizedLowerBound.DisjointCoordinate}
\label{CommunicationComplexity.Functions.Disjointness.RandomizedLowerBound.DisjointCoordinate}
\leanok
The inductive type with three constructors \texttt{neither}, \texttt{leftOnly},
\texttt{rightOnly}, encoding the three disjoint coordinate pairs
$(A_i, B_i) \in \{(0,0), (1,0), (0,1)\}$ placed at every coordinate $i$ other than the
special coordinate $T$ in the hard distribution for disjointness.
\end{definition}

\begin{theorem}[Cardinality of a singleton subtype]
\lean{CommunicationComplexity.Functions.Disjointness.RandomizedLowerBound.card_subtype_singleton_eq_one}
\label{CommunicationComplexity.Functions.Disjointness.RandomizedLowerBound.card_subtype_singleton_eq_one}
\leanok
For any type $\alpha$ and element $a \in \alpha$, the subtype $\{x : \alpha \mid x \in \{a\}\}$
has cardinality exactly $1$, for any \texttt{Fintype} instance on that subtype.
\end{theorem}
